\section{Du dividende brut \`a ton compte en Allemagne}\label{sec:fiscalite}
% Chapitre 3 (tache 10, passage chapitres 02-03). Figures: fig_steuer_100,
% fig_steuer_kumuliert (voir docs/figures.md). Aucun chiffre tape: macros de
% data/kennzahlen.tex. Chaque cle de data/quellen.json n'est citee qu'une fois
% dans ce chapitre (pas de note en double sur une meme page imprimee,
% scripts/check_footnote_pages.py). Titres d'axes et etiquettes passes par
% l'argument optionnel de \linienfigur/\balkenfigur (correctif infrastructure,
% tache 10, second passage), comme au chapitre 2.

Ce chapitre r\'epond \`a la question~: sur cent euros de dividende brut, combien arrivent vraiment sur
ton compte en Allemagne, selon le pays qui verse le dividende, selon que tu d\'etiens une action ou un
ETF, et que reste-t-il une fois l'assurance maladie pay\'ee~? Il d\'etaille les postes de la cascade du
chapitre~\ref{sec:depart}~: l'imp\^ot allemand (section~\ref{sec:fisc-abgeltung}), l'avantage des ETF
(section~\ref{sec:fisc-teilfreistellung}), les retenues \'etrang\`eres
(section~\ref{sec:fisc-quellensteuer}), l'imp\^ot d'\'Eglise (section~\ref{sec:fisc-kirchensteuer}),
l'option pour le bar\`eme progressif (section~\ref{sec:fisc-guenstiger}) et l'assurance maladie
(section~\ref{sec:fisc-kv}), puis chiffre le co\^ut fiscal cumul\'e de chaque placement
(section~\ref{sec:fisc-kumuliert}).

\bk{Ce chapitre d\'ecrit les r\`egles de 2026 telles que le mod\`ele les applique. Chaque point est \`a
faire valider par un conseiller fiscal avant d'engager de l'argent~; ce document n'est pas un conseil
fiscal.}

\annahme{R\`egles fiscales et sociales de 2026 fig\'ees pendant tout le plan~: les taux et le
forfait d'\'epargnant restent \`a leur niveau nominal de 2026, que la loi fixe en euros. En euros
constants, le forfait d'\'epargnant perd donc chaque ann\'ee un peu de sa valeur, ce que le mod\`ele
reproduit. Le plancher et le plafond de l'assiette de l'assurance maladie et d\'ependance suivent une
autre r\`egle~: la loi les revalorise chaque ann\'ee avec les salaires\QH{kv_schwellen_dynamisierung},
et le mod\`ele les tient donc constants en euros de 2026, ce qui revient \`a supposer que les salaires
progressent au moins comme les prix. Aucune r\'eforme future n'est anticip\'ee.}

\subsection{Abgeltungsteuer, Soli et Sparerpauschbetrag}\label{sec:fisc-abgeltung}

En Allemagne, les dividendes d'un particulier sont impos\'es \`a part de ses autres revenus, au taux
forfaitaire de \AbgeltungSatz~\% (Abgeltungsteuer)\QH{abgeltungsteuer_satz}, auquel s'ajoute la
contribution de solidarit\'e (Solidarit\"atszuschlag) de \SoliSatz~\% de cet imp\^ot\QH{soli_satz}. Le
pr\'el\`evement total atteint ainsi \SteuerSatzGesamt~\% du dividende. Chaque ann\'ee, les premiers
\Sparerpauschbetrag~€ de revenus de capitaux d'une personne seule restent exon\'er\'es
(Sparerpauschbetrag)\QH{sparerpauschbetrag}. Ce forfait compte pendant les premi\`eres ann\'ees
d'\'epargne, o\`u il couvre une bonne partie des dividendes, mais il ne p\`ese presque rien face aux
\BruttoNoetig~€ bruts annuels de la retraite vis\'ee. Au-del\`a du forfait, cent euros de dividende
allemand te laissent \NettoDe~€.

\subsection{Teilfreistellung~: pourquoi l'ETF est moins impos\'e}\label{sec:fisc-teilfreistellung}

Un fonds d'actions, et donc un ETF actions, b\'en\'eficie d'une exon\'eration partielle
(Teilfreistellung)~: \Teilfreistellung~\% de ses distributions et de ses plus-values sont exon\'er\'es
pour un particulier\QH{teilfreistellung_aktienfonds}. Cette r\`egle compense l'imp\^ot que le fonds
paie d\'ej\`a \`a l'int\'erieur, et les retenues \'etrang\`eres subies par le fonds sont d\'ej\`a
d\'eduites du rendement qu'il distribue. Sur cent euros distribu\'es par un ETF actions, il te reste
\NettoEtf~€. Un ETF capitalisant ne distribue rien, mais il est impos\'e chaque ann\'ee sur un montant
forfaitaire (Vorabpauschale) calcul\'e \`a partir d'un taux de r\'ef\'erence publi\'e par le
minist\`ere des Finances, le Basiszins, de \Basiszins~\% pour 2026\QH{basiszins_2026}~; la m\^eme
exon\'eration partielle s'y applique.

La cons\'equence est d\'efavorable \`a la th\`ese de ce plan, et il faut la dire d\`es maintenant~: \`a
dividende \'egal, une action d\'etenue en direct est plus lourdement impos\'ee en Allemagne qu'un ETF
d'actions. L'ETF maison du chapitre~\ref{sec:portefeuille} est un ensemble d'actions en direct, sans exon\'eration
partielle.

\subsection{Retenues \`a la source par pays et formulaire W-8BEN}\label{sec:fisc-quellensteuer}

Quand la soci\'et\'e est \'etrang\`ere, son pays retient d'abord un imp\^ot \`a la source.
L'Allemagne l'impute ensuite sur son propre imp\^ot, jusqu'au taux pr\'evu par la convention fiscale
entre les deux pays~; ce qui d\'epasse ce taux est perdu, sauf demande de remboursement aupr\`es du pays
d'origine\QH{quellensteuer}. Les \'Etats-Unis retiennent \QstUsStandard~\% sans formalit\'e, mais
seulement \QstUsFormulaire~\% si ton courtier dispose du formulaire W-8BEN, et ces \QstUsFormulaire~\%
sont enti\`erement imputables.

\annahme{\HypotheseUsFormulaire. C'est un \'ecart assum\'e au plan, qui partait du taux sans
formulaire~: le d\'ep\^ot est une d\'emarche unique, \`a v\'erifier \`a l'ouverture du compte, et sans
lui le d\'ecompte ci-dessous serait nettement plus mauvais pour les dividendes am\'ericains. La Suisse
et la France gardent leur retenue standard, sans remboursement.}

Avec le formulaire, cent euros de dividende am\'ericain te laissent \NettoUs~€, un peu plus qu'un
dividende allemand~: la retenue am\'ericaine r\'eduit l'imp\^ot allemand, et la contribution de
solidarit\'e ne porte que sur ce qui reste d'imp\^ot allemand. Sans formulaire, il ne reste que
\NettoUsSansFormulaire~€. La Suisse retient \QstCh~\%, dont seuls \QstChAnrechenbar~\% sont imputables
en Allemagne~: sans remboursement demand\'e en Suisse, il te reste \NettoCh~€, et \NettoChRemboursement~€
avec. La France retient \QstFr~\%, d\'ej\`a sous le taux de la convention et enti\`erement imputable,
d'o\`u \NettoFr~€ avec ou sans formulaire. Les Pays-Bas retiennent \QstNl~\% (\NettoNl~€ nets) et le
Royaume-Uni \QstGb~\% (\NettoGb~€, comme un dividende allemand). La figure~\ref{fig:net-cent-euros}
r\'esume ces cas.

% Ordre des etiquettes = ordre des lignes de data/fig_steuer_100.csv (DE, US,
% US_sans_w8ben, CH, CH_avec_remboursement, FR, FR_avec_formulaire, NL, GB, ETF,
% DE_kirchensteuer_8, DE_kirchensteuer_9).
\balkenfigur{fig_steuer_100}[ylabel={Net pour cent euros bruts (€)},
  ymin=0, enlarge y limits={upper, value=0.25}, bar width=11pt,
  xticklabels={{Allemagne},{\'Etats-Unis, W-8BEN},{\'Etats-Unis sans W-8BEN},{Suisse},
    {Suisse, rembours\'ee},{France},{France, formulaire},{Pays-Bas},{Royaume-Uni},{ETF actions},
    {Allemagne, \'Eglise \KistAcht~\%},{Allemagne, \'Eglise \KistNeun~\%}},
  nodes near coords, every node near coord/.append style={font=\tiny, rotate=90, anchor=west,
    /pgf/number format/fixed, /pgf/number format/fixed zerofill, /pgf/number format/precision=2}]{art}{netto}%
  {Montant net qui te reste pour cent euros de dividende brut, selon le pays de la soci\'et\'e ou
  l'enveloppe, au-del\`a du forfait d'\'epargnant (r\`egles de 2026, sans assurance
  maladie).\label{fig:net-cent-euros}}

Ce graphique montre que l'ETF actions laisse le plus, \NettoEtf~€ pour cent euros bruts, alors que les
actions en direct laissent entre \NettoCh~€ (Suisse sans remboursement) et \NettoUs~€ (\'Etats-Unis
avec formulaire, Pays-Bas). Les deux barres les plus coûteuses sont des oublis administratifs~: le
formulaire W-8BEN non d\'epos\'e (\NettoUsSansFormulaire~€) et le remboursement suisse non demand\'e.
Il faut toutefois comparer ce qui est comparable~: l'ETF monde distribue environ \RenditeDivEtf~\% de sa
valeur par an, contre \RenditeDivMaison~\% pour le portefeuille-exemple, et le chapitre~\ref{sec:capital} en tire le
capital n\'ecessaire pour chaque strat\'egie.

\subsection{Imp\^ot d'\'Eglise~: une variante}\label{sec:fisc-kirchensteuer}

\annahme{Pas d'imp\^ot d'\'Eglise (Kirchensteuer) dans le sc\'enario de base~; tu n'es pas suppos\'e
membre d'une \'Eglise qui le pr\'el\`eve. La variante le chiffre \`a \KistAcht~\% et \`a \KistNeun~\%
de l'imp\^ot, les deux taux possibles selon le Land de r\'esidence, \`a confirmer par toi.}

Si tu le paies, l'imp\^ot d'\'Eglise s'ajoute \`a l'Abgeltungsteuer, qui baisse l\'eg\`erement en
contrepartie, car la loi tient compte de sa d\'eductibilit\'e dans la formule du taux forfaitaire.
Cent euros de dividende allemand te laissent alors \NettoDeKistAcht~€ au taux de \KistAcht~\% et
\NettoDeKistNeun~€ au taux de \KistNeun~\%, au lieu de \NettoDe~€. \`A la retraite, le dividende brut
annuel n\'ecessaire passe de \BruttoNoetig~€ \`a \BruttoNoetigKistAcht~€ ou \BruttoNoetigKistNeun~€~:
l'\'ecart reste faible devant les autres postes de la cascade.

\subsection{G\"unstigerpr\"ufung~: le bar\`eme progressif sur demande}\label{sec:fisc-guenstiger}

Sur demande dans ta d\'eclaration de revenus, le fisc compare l'imp\^ot forfaitaire avec l'imp\^ot
que donnerait le bar\`eme progressif ordinaire sur tes revenus de capitaux, et retient le plus faible
(G\"unstigerpr\"ufung, \S~32d EStG, d\'ej\`a cit\'e plus haut). Pendant ta vie active, ton salaire
place tes revenus dans une tranche o\`u cette option ne sert \`a rien. \`A la retraite, si tu vis de
tes seuls dividendes, la situation change~: le bar\`eme s'applique \`a tout le dividende, tes
cotisations d'assurance maladie et d\'ependance deviennent d\'eductibles, et l'imp\^ot restant peut
tomber sous le seuil de la contribution de solidarit\'e.

\annahme{Calcul simplifi\'e, fait uniquement pour la cascade du chapitre~\ref{sec:depart}~: aucun
autre revenu, bar\`eme 2026 de l'imp\^ot sur le revenu\QH{est_tarif_2026}, cotisations maladie et
d\'ependance enti\`erement d\'eductibles, retenues \'etrang\`eres imput\'ees pays par pays dans la
limite de l'imp\^ot allemand correspondant, contribution de solidarit\'e nulle tant que l'imp\^ot reste
sous \SoliFreigrenze~€\QH{soli_freigrenze_2026}. Ce point est \`a faire valider par un conseiller.}

Pour les \BruttoNoetig~€ bruts de la cascade, le revenu imposable au bar\`eme tombe \`a
\GuenstigerRevenuImposable~€, et l'imp\^ot allemand apr\`es imputation des retenues \'etrang\`eres
n'est plus que de \GuenstigerImpotTarif~€, contre \ImpotAllemand~€ au taux forfaitaire. L'option vaut
donc environ \GuenstigerEconomieAn~€ par an, soit \GuenstigerEconomieMois~€ par mois. Le reste du
document garde volontairement le taux forfaitaire~: les capitaux des chapitres~\ref{sec:capital} et~\ref{sec:accumulation} sont donc
prudents sur ce point. \bk{Ne compte pas sur ce gain pour dimensionner ton plan~: il dispara\^it en
partie d\`es que ta rente l\'egale s'ajoute \`a tes revenus imposables, et il d\'epend d'un bar\`eme
qui aura chang\'e d'ici ta retraite.}

\subsection{Assurance maladie et d\'ependance volontaire}\label{sec:fisc-kv}

\annahme{Avant la rente l\'egale, tu restes affili\'e \`a l'assurance maladie publique (GKV) comme
membre volontaire, sans indemnit\'es journali\`eres (taux r\'eduit), sans enfant et avec la
cotisation compl\'ementaire moyenne de 2026. L'acc\`es \`a l'assurance des retrait\'es n'est pas
suppos\'e acquis.}

Un membre volontaire cotise sur l'ensemble de ses revenus, dividendes compris\QH{kv_teilfreistellung_etf}.
Le taux r\'eduit est de \KvSatz~\%\QH{kv_satz_ermaessigt}, la cotisation compl\'ementaire moyenne de
\KvZusatz~\%\QH{kv_zusatzbeitrag_2026}, et l'assurance d\'ependance d'une personne sans enfant de
\PvSatz~\%\QH{pv_satz_kinderlos_2026}, soit \KvPvSatz~\% au total. L'assiette est born\'ee des deux
c\^ot\'es. En bas, elle ne descend jamais sous \KvMindestMonat~€ par
mois\QH{kv_mindestbemessung_monat_2026}, ce qui fait une cotisation minimale de \KvPvMinMonat~€ par
mois m\^eme avec peu de revenus. En haut, elle s'arr\^ete \`a \KvBbgMonat~€ par
mois\QH{kv_bbg_monat_2026}, ce qui plafonne la cotisation \`a \KvPvMaxMonat~€ par mois. Comme ces deux
seuils sont revaloris\'es chaque ann\'ee (hypoth\`ese en t\^ete de chapitre), la cotisation maximale
reste de \KvPvMaxMonat~€ par mois en euros de 2026 \`a toutes les dates du plan~: l'inflation ne
l'all\`ege pas.

Avec \BruttoNoetigMois~€ de dividendes bruts par mois, tu d\'epasses ce plafond~: la cotisation est
au maximum, soit \KvPvAnnuel~€ par an, et chaque euro de dividende au-del\`a du plafond n'en supporte
plus. Pour les revenus de fonds, l'exon\'eration partielle des ETF s'applique aussi \`a l'assiette de
la cotisation, alors que les dividendes d'actions en direct y entrent en entier~: c'est le deuxi\`eme
avantage de l'ETF, dans la m\^eme source que ci-dessus.

\subsection{Co\^ut fiscal cumul\'e sur dix-huit ans}\label{sec:fisc-kumuliert}

Les sections pr\'ec\'edentes comparent les placements sur un seul dividende. La
figure~\ref{fig:impot-cumule-strategies} les compare sur toute une phase d'\'epargne.

\annahme{\'Epargne de \SparquoteZwanzig~\% du salaire net, \KapitalDepartMitte~€ de d\'epart, de 2027
\`a 2044, avec le rendement r\'eel du chapitre~\ref{sec:depart}. Ce taux commun sert seulement \`a
comparer les placements entre eux~; ce n'est pas celui du sc\'enario de base, qui \'epargne
\ReferenzSparquote~\%. Trois placements compar\'es~: tout en actions en direct (ETF maison seul), ETF
monde capitalisant (la r\'ef\'erence) et ETF monde distribuant. \HypotheseUmschichtung.}

\linienfigur{fig_steuer_kumuliert}[xlabel={Ann\'ee}, ylabel={Imp\^ots cumul\'es (euros de 2026)},
  xticklabel style={/pgf/number format/set thousands separator={}}, scaled y ticks=false,
  legend pos=north west]%
  {jahr}{maison,etf_thes,etf_aus}%
  {{Actions en direct},{ETF capitalisant},{ETF distribuant}}%
  {Imp\^ots cumul\'es pendant l'accumulation, retenues \'etrang\`eres comprises, selon le placement
  (\SparquoteZwanzig~\% d'\'epargne, euros de 2026).\label{fig:impot-cumule-strategies}}

Ce graphique montre qu'apr\`es dix-huit ans d'\'epargne, les actions en direct ont co\^ut\'e
\SteuerKumMaison~€ d'imp\^ots, retenues \'etrang\`eres comprises, contre \SteuerKumEtfThes~€ pour l'ETF capitalisant et \SteuerKumEtfAus~€
pour l'ETF distribuant~; l'\'ecart entre les actions en direct et la r\'ef\'erence atteint
\SteuerDifferenzAchtzehn~€. Les courbes des ETF restent \`a z\'ero les premi\`eres ann\'ees, parce que
leur montant imposable tient dans le forfait d'\'epargnant, alors que les actions en direct le
d\'epassent vite en distribuant \RenditeDivMaison~\% par an sans exon\'eration partielle.

Une limite du mod\`ele joue en faveur de l'ETF distribuant~: il ne calcule pas le compl\'ement de
Vorabpauschale d\^u quand la distribution, ici \RenditeDivEtf~\%, reste sous le rendement de base
d\'eriv\'e du Basiszins. L'\'ecart entre les deux courbes d'ETF ne doit donc pas \^etre lu comme un
avantage du distribuant. L'\'ecart avec les actions en direct, lui, est bien r\'eel, et il ne mesure
qu'une partie du co\^ut~: chaque euro d'imp\^ot pay\'e pendant l'\'epargne cesse de capitaliser, ce
que les projections des chapitres~\ref{sec:capital} et~\ref{sec:accumulation} prennent en compte.

\bk{En r\'esum\'e, la fiscalit\'e allemande favorise l'ETF sur trois points (exon\'eration partielle
de l'imp\^ot, exon\'eration partielle de l'assiette maladie, report de l'imp\^ot pour le capitalisant)
et p\'enalise deux oublis faciles \`a \'eviter (formulaire W-8BEN et remboursement suisse).} Le
chapitre~\ref{sec:capital} part de ces montants nets pour calculer le capital n\'ecessaire \`a chaque strat\'egie.

\Yannbox{\AnneeYannBasis}{Avant sa retraite, un regard en arri\`ere sur ses dix-huit premi\`eres
ann\'ees d'\'epargne montre ce que l'imp\^ot lui a d\'ej\`a co\^ut\'e en chemin~: entre 2027 et 2044, s'il n'avait d\'etenu que des actions en direct en
\'epargnant \SparquoteZwanzig~\% de son salaire, Yann aurait d\'ej\`a vers\'e \SteuerKumMaison~€ d'imp\^ots, retenues \'etrang\`eres comprises (euros de 2026), contre \SteuerKumEtfThes~€ avec
un ETF capitalisant. \`A la retraite, si tout son revenu venait d'actions en direct, chaque ann\'ee de
\BruttoNoetig~€ de dividendes bruts lui co\^uterait, aux r\`egles de 2026, \RetenuesEtrangeres~€ de
retenues \'etrang\`eres, \ImpotAllemand~€ d'imp\^ot allemand et \KvPvAnnuel~€ de cotisations maladie
et d\'ependance. En demandant la G\"unstigerpr\"ufung, il garderait environ \GuenstigerEconomieMois~€
de plus par mois.}
